% Bagean: sajarah
% Bab: biografi
% Pérangan: rozsa-peter

\documentclass[../../../include/open-logic-section]{subfiles}

\begin{document}

\olfileid{his}{bio}{pet}

\olsection{R\'ozsa P\'eter}

\olphoto{peter-rozsa}{R\'ozsa P\'eter}

R\'ozsa P\'eter lair kanthi jeneng
R\'ozsa Politzer ing Budapest, Hongaria,
tanggal 17 Februari 1905. Dhèwèké
paling misuwur amarga pakaryané ngenani
fungsi rekursif, kang wigati banget kanggo
lairé bidhang téori rekursi.

P\'eter gedhé ing mangsa pulitik kang
abot---Perang Donya I dumadi nalika
dhèwèké isih remaja---nanging bisa
sekolah ing Sekolah Putri Maria
Terezia kang makmur ing Budapest,
lan lulus saka kono ing taun 1922.
Banjur dhèwèké kuliah ing
Universitas P\'azm\'any P\'eter
(kang banjur dijenengi Universitas
Lor\'and E\"otv\"os) ing Budapest.
Miturut karepé bapaké, wiwitané
dhèwèké sinau kimia, nanging
banjur pindhah menyang matématika
lan lulus ing taun 1927.
Senajan nduwèni kualifikasi
kanggo mulang matématika ing
sekolah menengah, kahanan ékonomi
wektu iku abot banget amarga
Depresi Gedhé mengaruhi
ékonomi donya. Ing mangsa iki,
P\'eter nampa warna-warna
pakaryan minangka tutor lan
guru matématika pribadi.
Pungkasané dhèwèké bali
menyang universitas kanggo
sinau matématika ing tingkat
lanjut. Wiwitané dhèwèké
ngrancang nliti téori bilangan,
nanging sawisé mangertèni
yèn asilé wis dibuktèkaké
déning wong liya, dhèwèké
meh ninggalaké matématika
babar pisan. Dhèwèké
didorong kanggo nggarap
teorema-teorema
ketaklengkapan G\"odel,
lan tanpa sadhar mbuktèkaké
sawetara asilé kanthi cara
kang béda. Bab iki mbalèkaké
kapercayan dhiriné, lan
P\'eter banjur nulis
makalah-makalah kapisané
ngenani téori rekursi,
kanthi inspirasi saka
program dhasar matématika
David Hilbert. Dhèwèké
nampa gelar doktor ing taun
1935, lan ing taun 1937
dadi rédaktur \emph{Journal of
  Symbolic Logic}.

Makalah-makalah wiwitan
P\'eter dianggep sacara
amba minangka sumbangan
pangadeg kanggo téori
fungsi rekursif. Ing
\cite{Peter1935a}, dhèwèké
nliti gegayutan antarane
warna-warna rekursi.
Ing \cite{Peter1935b},
dhèwèké nuduhaké yèn
sawijining fungsi kang
didhéfinisikaké kanthi
rekursif dudu fungsi
rekursif primitif. Iki
nyederhanakaké asil
sadurungé saka Wilhelm
Ackermann. Fungsi kang
disederhanakaké déning
P\'eter iki saiki kerep
diarani fungsi Ackermann---
lan kadhangkala, luwih
trep, fungsi Ackermann--P\'eter.
Dhèwèké nulis buku
kapisan ngenani téori
fungsi rekursif
\citep{Peter1951}.

Senajan pakaryané wigati
lan duwé pangaruh,
P\'eter durung olèh
jabatan mulang tetep
nganti taun 1945.
Nalika Nazi ndhudhuki
Hongaria sajroning
Perang Donya II,
P\'eter ora diidinaké
mulang amarga
undhang-undhang
anti-Semit. Ing taun
1944 pamaréntah
ngedegaké ghetto
Yahudi ing Budapest;
ghetto iku kapisah
saka kutha liyané
lan dijaga déning
penjaga mawa gaman.
P\'eter dipeksa
manggon ing ghetto
nganti dibébaské
ing taun 1945.
Banjur dhèwèké
mulang ing
Budapest Teachers
Training College,
lan wiwit taun
1955 ing Universitas
E\"otv\"os Lor\'and.
Dhèwèké matématikawan
wanita Hongaria
kapisan kang dadi
Doktor Akadhemik
Matématika, lan
wanita kapisan
kang kapilih dadi
anggota Akademi
Ilmu Pengetahuan
Hongaria.

P\'{e}ter kawentar
minangka guru
matématika kang
kebak semangat,
lan luwih seneng
nyinaoni sipat
lan kaéndahan
masalah matématika
bareng murid-muridé
tinimbang mung
mènèhi kuliah.
Amarga iku,
murid-muridé
nyebut dhèwèké
kanthi tresna
``Bulik Rosa''.
P\'eter séda
ing taun 1977
nalika umur~71 taun.

\begin{reading}
Kanggo wacan biografi
liyane, delengen
\citep{Oconnor2014}
lan \citep{Andrasfai1986}.
\citet{Tamassy1994}
nindakaké wawancara
ringkes karo P\'eter.
Kanggo wacan matématika
kang nyenengaké,
delengen buku P\'eter
\emph{Playing With Infinity}
\citep{Peter2010}.
\end{reading}

\end{document}
